% reserves.tex — Reserve curves and coordinate transform verification
% Kontrol: prove_cfmm_reserve*

\section{Reserve Analysis}\label{sec:reserves}

This section verifies the reserve curves derived in \cref{sec:trading-function}
and establishes boundary behavior and numerical properties.

\subsection{Reserve Curves in Price Space}

Per unit liquidity ($L = 1$), with $p \in [\pl, \pu]$:

\begin{center}
\begin{tabular}{@{}lccc@{}}
\toprule
& $p = \pl$ & $p \in (\pl, \pu)$ & $p = \pu$ \\
\midrule
$x(p) = 2p/\pl^2$ & $2/\pl$ & $2p/\pl^2$ & $2\pu/\pl^2$ \\
$u(p) = x^2 = 4p^2/\pl^4$ & $4/\pl^2$ & $4p^2/\pl^4$ & $4\pu^2/\pl^4$ \\
$y(p) = (\pu^2+p^2)/\pl^2$ & $(\pu^2+\pl^2)/\pl^2$ & $(\pu^2+p^2)/\pl^2$ & $2\pu^2/\pl^2$ \\
$V(p) = (\pu^2-p^2)/\pl^2$ & $(\pu^2-\pl^2)/\pl^2$ & $(\pu^2-p^2)/\pl^2$ & $0$ \\
\bottomrule
\end{tabular}
\end{center}

\subsection{Monotonicity}

\begin{proposition}[Reserve Monotonicity]\label{prop:reserve-monotone}
Both $x(p)$ and $y(p)$ are strictly increasing in $p$ on $[\pl, \pu]$:
\[
  x'(p) = \frac{2}{\pl^2} > 0, \qquad
  y'(p) = \frac{2p}{\pl^2} > 0 \quad\text{for } p > 0
\]
\end{proposition}

\begin{remark}
Both reserves increasing with $p$ is the correct direction under Framing~C
(\cref{sec:architecture}): as concentration rises, more virtual exposure
is absorbed ($x$ rises) and more numeraire accumulates ($y$ rises) because
the underwriter's collateral drains into the payoff liability. The LP value
$V(p) = y - px$ decreases, which is the underwriter bearing losses.
\end{remark}

\subsection{Collateral Accounting}

The underwriter deposits USDC (Token~Y). The virtual Token~X reserve is
an internal accounting variable with no external token.

\begin{definition}[Collateral Value]\label{def:collateral-value}
At price $p$, the underwriter's real (USDC) collateral per unit liquidity equals
the LP value:
\begin{equation}\label{eq:collateral-value}
  \text{collateral}(p) = V(p) = y(p) - p \cdot x(p) = \frac{\pu^2 - p^2}{\pl^2}
\end{equation}
At $p = \pl$: full collateral $(\pu^2 - \pl^2)/\pl^2$.
At $p = \pu$: zero (fully drained). The underwriter cannot lose more than
their initial deposit---the CLAMM tick cap $\pu$ enforces bounded loss.
\end{definition}

\subsection{Numerical Verification}

For $\DeltaStar = 0.09$, $\pl = 0.0989$, $\pu = 1.0$ ($\DeltaPlus_{\max} = 0.5$):

\begin{center}
\begin{tabular}{@{}lcccc@{}}
\toprule
$p$ & $\DeltaPlus$ & $x$ & $y$ & $V$ (USDC/unit $L$) \\
\midrule
$0.0989$ ($= \pl$) & $0.09$ & $20.22$ & $103.2$ & $101.2$ \\
$0.1758$ & $0.149$ & $35.95$ & $105.4$ & $99.04$ \\
$0.5$ & $0.333$ & $102.2$ & $127.6$ & $76.5$ \\
$1.0$ ($= \pu$) & $0.5$ & $204.5$ & $204.5$ & $0$ \\
\bottomrule
\end{tabular}
\end{center}

Payoff check at $p = 0.1758$: $N = V(\pl) - V(p) = 101.2 - 99.04 = 2.16$.
Formula: $(0.1758/0.0989)^2 - 1 = 3.16 - 1 = 2.16$ \checkmark.
